\documentclass[tikz,border=10pt]{standalone}
\usepackage[utf8]{inputenc}
\usepackage{ctex}
\usepackage{tikz}
\usetikzlibrary{shapes,arrows,positioning,fit,calc,decorations.pathreplacing}

\begin{document}

% Transformer 解码器架构图
\begin{tikzpicture}[
    % 节点样式
    block/.style={rectangle, draw, fill=blue!10, rounded corners=4pt, 
                  minimum width=3cm, minimum height=0.8cm, align=center,
                  font=\small\sffamily, inner sep=4pt},
    emb/.style={rectangle, draw, fill=orange!20, rounded corners=4pt,
                minimum width=3cm, minimum height=0.7cm, align=center,
                font=\small\sffamily},
    attn/.style={rectangle, draw, fill=green!15, rounded corners=4pt,
                 minimum width=3.2cm, minimum height=0.9cm, align=center,
                 font=\small\sffamily, inner sep=4pt},
    ffn/.style={rectangle, draw, fill=purple!10, rounded corners=4pt,
                minimum width=3cm, minimum height=0.8cm, align=center,
                font=\small\sffamily},
    addnorm/.style={rectangle, draw=gray!60, dashed, fill=gray!5, rounded corners=3pt,
                    minimum width=2.6cm, minimum height=0.5cm, align=center,
                    font=\scriptsize\sffamily},
    arrow/.style={->, >=stealth, thick, color=black!70},
    brace_arrow/.style={decorate, decoration={brace, amplitude=6pt, mirror}, 
                        thick, black!60},
    label_font/.style={font=\scriptsize\sffamily, color=gray!70},
    % 颜色
    data_color/.style={fill=cyan!20, draw=cyan!60!black},
    out_color/.style={fill=red!15, draw=red!60!black},
]

% ===== 左半部分：编码器 (Encoder) =====
\begin{scope}[shift={(-5.5,0)}]
    % 标题
    \node[font=\large\sffamily\bfseries] at (0, 9.3) {编码器 (Encoder)};
    
    % 输入嵌入
    \node[emb, data_color] (enc_in) at (0, 8.2) {输入嵌入};
    
    % 位置编码
    \node[emb, fill=yellow!15, draw=yellow!60!black] (enc_pos) at (0, 7.2) {位置编码};
    \draw[arrow] (enc_in) -- (enc_pos);
    
    % 多头自注意力
    \node[attn] (enc_attn) at (0, 5.6) {多头自注意力\\ (Masked Self-Attention)};
    \draw[arrow] (enc_pos) -- (enc_attn);
    
    % 残差+层归一化
    \node[addnorm] (enc_norm1) at (0, 4.5) {Add \& LayerNorm};
    \draw[arrow] (enc_attn) -- (enc_norm1);
    
    % 残差连接
    \draw[arrow, gray!50] ($(enc_pos.east)+(0.2,0)$) -- 
          ($(enc_pos.east)+(0.8,0)$) |- ($(enc_norm1.east)+(0.5,0)$);
    \node[label_font] at ($(enc_pos.east)+(1.8,0.8)$) {残差连接};
    
    % 前馈网络
    \node[ffn] (enc_ffn) at (0, 3.0) {前馈网络 (FFN)};
    \draw[arrow] (enc_norm1) -- (enc_ffn);
    
    % 残差+层归一化
    \node[addnorm] (enc_norm2) at (0, 1.9) {Add \& LayerNorm};
    \draw[arrow] (enc_ffn) -- (enc_norm2);
    
    % 残差连接
    \draw[arrow, gray!50] ($(enc_norm1.east)+(0.2,0)$) -- 
          ($(enc_norm1.east)+(0.8,0)$) |- ($(enc_norm2.east)+(0.5,0)$);
    
    % 编码器输出
    \node[block, fill=teal!10, draw=teal!60!black] (enc_out) at (0, 0.7) {编码器输出};
    \draw[arrow] (enc_norm2) -- (enc_out);
    
    % 重复标记
    \node[label_font, font=\tiny\sffamily] at (-2.2, 6.4) {$\times N$};
    \draw[gray!30] ($(-1.6,6.9)$) rectangle ($(-1.4,5.1)$);
    
    % 侧边标注
    \draw[brace_arrow] ($(enc_in.west)+(-0.5,0.35)$) -- 
          ($(enc_out.west)+(-0.5,-0.35)$);
    \node[label_font, font=\footnotesize\sffamily] at ($(enc_in.west)+(-2.0,0)$) 
          {编码器层 $\times 6$};
\end{scope}

% ===== 右半部分：解码器 (Decoder) =====
\begin{scope}[shift={(5.5,0)}]
    % 标题
    \node[font=\large\sffamily\bfseries] at (0, 10.6) {解码器 (Decoder)};
    
    % 输出嵌入（右移）
    \node[emb, data_color] (dec_in) at (0, 9.5) {输出嵌入（右移）};
    
    % 位置编码
    \node[emb, fill=yellow!15, draw=yellow!60!black] (dec_pos) at (0, 8.5) {位置编码};
    \draw[arrow] (dec_in) -- (dec_pos);
    
    % 掩码多头自注意力
    \node[attn, fill=green!20] (dec_attn1) at (0, 6.9) {掩码多头自注意力\\ (Masked Self-Attention)};
    \draw[arrow] (dec_pos) -- (dec_attn1);
    
    % Add & LayerNorm
    \node[addnorm] (dec_norm1) at (0, 5.8) {Add \& LayerNorm};
    \draw[arrow] (dec_attn1) -- (dec_norm1);
    
    % 残差连接
    \draw[arrow, gray!50] ($(dec_pos.east)+(0.2,0)$) -- 
          ($(dec_pos.east)+(0.8,0)$) |- ($(dec_norm1.east)+(0.5,0)$);
    
    % 交叉注意力
    \node[attn, fill=cyan!20, draw=cyan!60!black] (dec_cross) at (0, 4.2) 
          {交叉注意力\\ (Cross-Attention)};
    \draw[arrow] (dec_norm1) -- (dec_cross);
    
    % 来自编码器的 KV
    \draw[arrow, teal!70!black] ($(enc_out.east)+(0.3,0)$) -- 
          ($(enc_out.east)+(1.2,0)$) |- ($(dec_cross.west)+(-0.05,0.2)$);
    \node[label_font, teal!70!black] at ($(enc_out.east)+(2.2,0.15)$) {Key / Value};
    
    % Add & LayerNorm
    \node[addnorm] (dec_norm2) at (0, 3.1) {Add \& LayerNorm};
    \draw[arrow] (dec_cross) -- (dec_norm2);
    
    % 残差连接
    \draw[arrow, gray!50] ($(dec_norm1.east)+(0.2,0)$) -- 
          ($(dec_norm1.east)+(0.8,0)$) |- ($(dec_norm2.east)+(0.5,0)$);
    
    % 前馈网络
    \node[ffn] (dec_ffn) at (0, 1.6) {前馈网络 (FFN)};
    \draw[arrow] (dec_norm2) -- (dec_ffn);
    
    % Add & LayerNorm
    \node[addnorm] (dec_norm3) at (0, 0.5) {Add \& LayerNorm};
    \draw[arrow] (dec_ffn) -- (dec_norm3);
    
    % 残差连接
    \draw[arrow, gray!50] ($(dec_norm2.east)+(0.2,0)$) -- 
          ($(dec_norm2.east)+(0.8,0)$) |- ($(dec_norm3.east)+(0.5,0)$);
    
    % 线性层 + Softmax
    \node[block, out_color] (dec_out) at (0, -0.8) {线性层 + Softmax};
    \draw[arrow] (dec_norm3) -- (dec_out);
    
    % 输出
    \node[block, fill=red!20, draw=red!80!black, minimum height=0.6cm] 
          (output) at (0, -1.9) {输出概率分布};
    \draw[arrow] (dec_out) -- (output);
    
    % 重复标记
    \node[label_font, font=\tiny\sffamily] at (-2.2, 7.7) {$\times N$};
    \draw[gray!30] ($(-1.6,8.2)$) rectangle ($(-1.4,6.4)$);
    
    % 侧边标注
    \draw[brace_arrow] ($(dec_in.west)+(-0.5,0.35)$) -- 
          ($(dec_ffn.east)+(-0.5,0)+(0,-0.5)$);
    \node[label_font, font=\footnotesize\sffamily] at ($(dec_in.west)+(-2.1,1.0)$) 
          {解码器层 $\times 6$};
\end{scope}

% ===== 中间连接箭头 =====
\node[align=center, font=\small\sffamily] at (0, 0.2) {编码器-解码器\\ 交互};
\draw[->, >=stealth, thick, dashed, teal!60!black] 
      ($(enc_out.east)+(0.2,0)$) -- ($(enc_out.east)+(0.6,0)$);
\draw[->, >=stealth, thick, dashed, teal!60!black] 
      ($(dec_cross.south)+(1.5,0)$) -- ($(dec_cross.south)+(1.1,0)$);

% ===== 整体标题 =====
\node[font=\huge\sffamily\bfseries, align=center] at (0, 11.8) 
      {Transformer 模型架构\\[3pt]
       {\normalsize\sffamily\mdseries —— 基于 《Attention Is All You Need》}};

\end{tikzpicture}

\end{document}
